\section{Conclusion}
This report presented a structured analysis of the Festo factory network and its associated industrial communication environment. The assessment covered the existing inventory of operational technology assets, communication protocols, observed network traffic through packet captures, and the resulting security implications in an industrial control system context.

The inventory phase identified a heterogeneous mix of OT and IT components, including PLCs, HMIs, industrial controllers, MQTT-enabled devices, OPC UA servers, and supporting IT infrastructure such as MES systems and engineering workstations. This diversity of assets reflects a typical Industry 4.0 environment, where legacy industrial systems coexist with modern networked components, increasing both interoperability and security complexity.

Protocol analysis highlighted the use of industrial standards such as OPC UA and MQTT alongside HTTP-based services and supporting infrastructure protocols. While OPC UA provides a comparatively strong security model when properly configured, other protocols observed in the environment introduce varying levels of exposure depending on their configuration and deployment context. The network capture analysis further demonstrated how inter-zone communication patterns can lead to unnecessary exposure of OT assets if segmentation and filtering are not strictly enforced.

Based on the identified assets and traffic flows, a threat analysis was conducted to evaluate potential attack vectors, including unauthorized lateral movement between IT and OT zones, interception or manipulation of industrial messages, and abuse of weak or misconfigured services. These threats were systematically evaluated in a risk matrix, which highlighted that the highest risks are associated with insufficient network segmentation, weak authentication boundaries, and uncontrolled inter-zone communication paths.

To address these findings, a set of security requirements was defined, aligned primarily with IEC 62443, ISO/IEC 27001, and relevant sections of the NIS2 directive. These requirements formalise controls for network segmentation, firewall enforcement, secure protocol configuration, authentication and access control, logging and monitoring, time synchronisation, and backup and recovery strategies.

Finally, a proposed architecture was developed based on a zone-and-conduit model consistent with IEC 62443 principles. The design introduces clearly separated OT Field, OT Cell, Industrial DMZ, and IT zones, with strictly controlled communication paths enforced through industrial firewalls and dedicated proxy services. This segmentation significantly reduces the attack surface by ensuring that critical OT assets are not directly accessible from higher-level IT networks and that all inter-zone communication is explicitly mediated and logged.

Overall, the proposed security architecture improves the resilience, traceability, and controllability of the Festo factory network. While no industrial environment can be fully risk-free, the combination of structured segmentation, protocol hardening, and aligned security standards provides a robust baseline for mitigating the most relevant threats identified in this analysis.